\section{A three-quarter shape restriction for ordinary alternating cores}
\label{sec:ordinary-alternating-three-quarters}

\paragraph{Status.}
New root derivation, with an author whole-proof reading and an
independent whole-proof review by /root/g34\_eightfolds, both PASS
on 21 November 2026. The independent arithmetic reconstruction passed. The exact finite certificate checks the
formulas and endpoints; the proof below supplies the all-degree
statement. The published content and low-degree inputs are exactly
those already read and cited in HD-ALT-01--03.

\begin{lemma}[HD-ALT-04a: a Cauchy improvement of the content bound]
\label{HD-ALT-04a}
Let $n\ge15$ and let $\lambda\vdash n$ satisfy
\[
                k:=\max(\lambda_1,\lambda'_1)\le3n/4.
                                                               \tag{ALTQ.1}
\]
For either irreducible constituent $W$ of $S^\lambda|_{A_n}$,
of degree $a$, put $u=\chi((123))/a$ and
$v=\chi((12)(34))/a$. Then
\[
                u<9/16,\qquad v<9/16.
                                                               \tag{ALTQ.2}
\]
Every nonscalar finite-order linear normalizer of the represented
core satisfies $\nu(T)\ge7a/64$, with every scalar phase allowed.
\end{lemma}

\begin{proof}
Write $t=\sum_{(i,j)\in\lambda}(j-i)$ and
$S=\sum_{(i,j)\in\lambda}(j-i)^2$. The exact specializations
(ALT.6) and (ALTN.6), including their split-constituent justification,
give
\[
 u=\frac{3(S-\binom n2)}{n(n-1)(n-2)},\qquad
 v=\frac{4(t^2+n(n-1)-3S)}{n(n-1)(n-2)(n-3)}.
                                                               \tag{ALTQ.3}
\]
The diagonal-hook estimate in HD-ALT-01b holds for any upper
bound $k$ on both the first row and the first column:
\[
 S\le\frac{k(2k-1)(n-1)}6.
\]
Its proof uses only that every arm and leg has length at most
$k-1$ and that the sum of their lengths is at most $n-1$.
Thus, by (ALTQ.1),
\[
 u\le\frac{k(2k-1)-3n}{2n(n-2)}
    \le\frac{9n-30}{16(n-2)}<\frac9{16}.
                                                               \tag{ALTQ.4}
\]

The row/column estimate (ALTN.7), before its specialization of $k$,
and its transpose give
\[
 |t|\le\frac{n(k-1)}2-(n-k)
       \le\frac{3n(n-2)}8.
                                                               \tag{ALTQ.5}
\]
Apply Cauchy's inequality to the $n$ contents to obtain $S\ge t^2/n$.
Keeping this term, rather than discarding $-3S$ in (ALTQ.3), yields
\[
 \begin{split}
 v&\le\frac{4t^2}{n^2(n-1)(n-2)}
              +\frac4{(n-2)(n-3)}\\
  &\le\frac{9(n-2)}{16(n-1)}
              +\frac4{(n-2)(n-3)}<\frac9{16}.
 \end{split}                                                   \tag{ALTQ.6}
\]
The final strict inequality is equivalent to
$9n^2-109n+118>0$. Its value at $15$ is $508$,
and its forward difference is $18n-100>0$ for $n\ge15$.
This proves (ALTQ.2). The complete normalizer bound HD-ALT-03a
now gives
\[
 \nu(T)\ge a\min\{(1-u)/3,(1-v)/4\}\ge7a/64.
\]
No assertion about the reality of a split constituent is needed
for this lemma.
\end{proof}

\begin{theorem}[HD-ALT-04b: transitive ordinary cores require a longer shape]
\label{HD-ALT-04b}
Consider an inducing block of an actual rank-nine Weil sixfold
finite-cover source. Suppose its finite image $J$ has an irreducibly
acting normal central-product core with simple factors $A_n$,
$n\ge15$, and ordinary factor representations. Suppose that $J$
acts transitively on the factors and that its character field
contains the original imaginary quadratic field $K$.

Then every factor's Specht label satisfies
\[
                   \max(\lambda_1,\lambda'_1)>3n/4.
                                                               \tag{ALTQ.7}
\]
In particular this is a proper refinement of the remaining
ordinary alternating gate of D-BI-19 and HD-ALT-03.
It permits all scalar phases, compatible outer automorphisms
and tensor permutations. No additional nonsplit-polarization
or inducing-field-equality hypothesis is required.
\end{theorem}

\begin{proof}
Suppose (ALTQ.1) holds for one factor. Transitivity implies that
all factor labels agree up to transpose and the alternating
restriction ambiguity, so (ALTQ.1) holds for every factor.

Each factor has degree at least $n^3/27\ge125$.
Indeed the precise published consequence of Rasala's classification
in \cite[proof of Theorem~9.7, p.~312]{HD:GKT2019} says that any
nonlinear Specht character of degree below $2n^3/27$ has, up to
transpose, one of the three shapes $(n-1,1),(n-2,2),(n-2,1,1)$.
For $n\ge15$ each has a row of length at least $n-2>3n/4$.
They cannot occur here. Passing to an alternating constituent
loses a factor of at most two and proves the bound.

A self-conjugate factor label has first row and column at most
$\lceil n/2\rceil$. By transitivity all factors would then satisfy
the stronger shape hypothesis of HD-ALT-03c, which already
excludes these data. We can therefore assume that no factor
label is self-conjugate. Every factor is consequently the
irreducible restriction of a real Specht representation.
The tensor-product core has a unique invariant line of symmetric
nondegenerate bilinear forms. Every element of $J$ preserves
this line, so acts by a multiplier character
$\mu:J\longrightarrow\mathbf C^*$.

This character is nontrivial. Otherwise $J$ preserves the form,
its representation is self-dual, and its finite-group character
is real valued. Its character field would be a subfield of
$\mathbf R$, contrary to $K\subseteq\mathbf Q(\chi_J)$.

Let $e$ be the total inducing degree. The actual intermediate-cover
factorization used in D-BI-01 and HD-ALT-03c forces genus zero,
at least twelve genuine inertias, and total moving rank $2e+6$:
genus at least two is impossible for $e\ge125$, and genus one
has period-parameter dimension at most six. The genus-zero
inertias generate $J$ and have product one. Therefore at least
two have nontrivial multiplier. For either such inertia, its
fixed space is totally isotropic for the preserved form line,
so its ordinary moving rank is at least $\lceil e/2\rceil$.

Every other genuine inertia moves at least $7e/64$ dimensions.
For factor-preserving operators this follows from HD-ALT-04a
and the largest-eigenspace tensor bound. A nontrivial factor
permutation has the stronger quarter bound from the weighted
tensor-cycle argument in D-SPOR-03. A genuine scalar inertia
moves all $e$ dimensions. Thus the exact necessary budget implies
\[
 2\lceil e/2\rceil+10\lceil7e/64\rceil\le2e+6,
 \qquad e+70e/64\le2e+6,
 \qquad e\le64.
                                                               \tag{ALTQ.8}
\]
This contradicts $e\ge125$ and proves (ALTQ.7).
\end{proof}

\paragraph{Meaning and boundary.}
This excludes a strictly larger infinite family of ordinary
alternating local systems from dominant Weil sixfold suppliers.
For example the nonselfconjugate shape $(11,4)$ at $n=15$
has $11>\lceil15/2\rceil$ and $11\le45/4$; it lies outside
the earlier half-size region and inside (ALTQ.1).
The conclusion concerns the finite monodromy representation and
the actual period-map factorization. Shapes with a still longer
row or column remain open unless covered by the separate standard,
exterior-square or unordered-pair results. No new cycle is produced.
In the project's language this narrows the supplier gate for
condition A (cote~16, scan~050). It does not strengthen condition I
(scan~049 and cote~44, Theorem~4), change the Cayley--Hamilton
projector in scans~048--049, or supply a new assertion about
weights or motivic $L$-functions.
